% content/hindi_section.tex
% ── Hindi Section ─────────────────────────────────────────────
%
% Language: Hindi (Modern Standard Hindi in Devanagari)
% Font family: \hin  (Scale=.85, same Shobhika-Bold)
%
% Use \begin{hindi}...\end{hindi} for block passages, or
% \texthindi{...} for inline snippets.

\setchapterhead{[अनुभाग का नाम / Section Name]}
\onehalfspacing

\begin{center}
  {\large\hin [मुख्य शीर्षक / Main Title]}\\
  \sectionrule[.2\linewidth]
\end{center}

\begin{hindi}

\noindent\textbf{[उपशीर्षक / Subheading १]}

[यहाँ हिन्दी पाठ लिखें / Enter Hindi text here.]

\vspace{3mm}
\noindent\textbf{[उपशीर्षक / Subheading २]}

[यहाँ हिन्दी पाठ लिखें / Enter Hindi text here.]

\vspace{3mm}
\noindent\textbf{[उपशीर्षक / Subheading ३]}

[यहाँ हिन्दी पाठ लिखें / Enter Hindi text here.]

\end{hindi}

% ── ADD MORE HINDI PARAGRAPHS BELOW ──────────────────────────
% Always wrap Hindi text in \begin{hindi}...\end{hindi}
% or use \texthindi{<inline text>}.
%
% To mix Hindi gloss with Sanskrit verse:
%
%   \begin{shloka}
%     <Sanskrit verse in \la font>
%   \end{shloka}
%   \begin{hindi}
%     [उपर्युक्त श्लोक का अर्थ है ...]
%   \end{hindi}

\newpage